\subsection{PROJECTIVE COMPLETION OF AN AFFINE SPACE}

Let V be a (left) vector space over a field K and consider the vector space \(K_{s} \times V\) over K; the projective space \(\mathbf{P}(\mathbf{K}_{s} \times \mathbf{V})\) is called the projective space canonically associated with the vector space V. If V is of dimension n, \(\mathbf{P}(\mathbf{K}_{s} \times \mathbf{V})\) is of the same dimension n.~Consider in \(K_{s} \times V\) the affine hyperplane \(V_{1} = \{1\} \times V\) , whose direction (no. 3) is the subspace \(V_{0} = \{0\} \times V\) ; if a line (passing through 0) of \(K_{s} \times V\) is not contained in \(V_{0}\) , it contains a point \((\alpha, x)\) with \(\alpha \neq 0\) and \(x \in V\) , hence it contains also the point \(\alpha^{-1}(\alpha, x) = (1, \alpha^{-1}x)\) of \(V_{1}\) ; the converse is immediate and it is seen that there is a one-to-one correspondence between the points of \(V_{1}\) and the lines (passing through 0) of \(\mathbf{K}_s \times \mathbf{V}\) not contained in \(\mathbf{V}_0\) , each of the latter meeting \(\mathbf{V}_1\) in one and only one point. It follows that the mapping \(x \mapsto \phi(x) = \pi(1, x)\) is an injection (called canonical) of \(\mathbf{V}\) into the projective space \(\mathbf{P}(\mathbf{K}_s \times \mathbf{V})\) ; \(\mathbf{V}\) is often identified with its image under this injection. The complement of \(\phi(\mathbf{V})\) in \(\mathbf{P}(\mathbf{K}_s \times \mathbf{V})\) is the projective hyperplane \(\mathbf{P}(\mathbf{V}_0)\) called the hyperplane at infinity of \(\mathbf{P}(\mathbf{K}_s \times \mathbf{V})\) (or of \(\mathbf{V}\) , by an abuse of language); its points are also called the ``points at infinity'' of \(\mathbf{P}(\mathbf{K}_s \times \mathbf{V})\) (or of \(\mathbf{V}\) ). If \((a_t)\) is a basis of \(\mathbf{V}\) and in \(\mathbf{K}_s \times \mathbf{V}\) the basis is taken consisting of the elements \(e_t = (0, a_t)\) and the element \(e_\omega = (1, 0)\) , the points at infinity in \(\mathbf{P}(\mathbf{K}_s \times \mathbf{V})\) are those whose homogeneous coordinate of index \(\omega\) is 0.

Let M be an affine linear variety in V (no. 3) and D its direction; the canonical image \(\phi(M)\) of M in \(\mathbf{P}(\mathbf{K}_s \times \mathbf{V})\) is contained in the canonical image \(\overline{\mathbf{M}} = \pi(\mathbf{M}_2)\) of the vector subspace \(\mathbf{M}_2\) of \(\mathbf{K}_s \times \mathbf{V}\) generated by the affine variety \(\mathbf{M}_1 = \{1\} \times \mathbf{M}\) of \(\mathbf{K}_s \times \mathbf{V}\) . More precisely, if \((a_t)\) is an affinely free system of M generating M, the elements \((1, a_t)\) form a basis of \(\mathbf{M}_2\) and therefore \(\overline{\mathbf{M}}\) is just the projective linear variety generated by \(\phi(M)\) ; if M is finite-dimensional, \(\overline{\mathbf{M}}\) has the same dimension as M. The complement of \(\phi(M)\) in \(\overline{\mathbf{M}}\) is the intersection of \(\overline{\mathbf{M}}\) and the hyperplane at infinity and is equal to the canonical image \(\pi(\mathbf{M}_0)\) , where \(\mathbf{M}_0 = \{0\} \times \mathbf{D}\) .

Conversely, let N be a projective linear variety not contained in the hyperplane at infinity and let \(\mathbf{R} = \overline{\pi}^{-1}(\mathbf{N})\) ; \(R \cap V_{1}\) is an affine linear variety of \(K \times V\) of the form \(\{1\} \times M\) , where M is an affine linear variety of V and it is immediately seen that N is the affine linear variety \(\overline{M}\) generated by \(\phi(M)\) .

There is therefore a one-to-one correspondence between the affine linear varieties of V and the projective linear varieties of \(\mathbf{P}(\mathbf{K}_{s} \times \mathbf{V})\) not contained in the hyperplane at infinity; for two affine linear varieties of V to be parallel, it is necessary and sufficient that the projective linear varieties which they generate have the same intersection with the hyperplane at infinity (which is sometimes expressed by saying that the two affine linear varieties in question have the same points at infinity).

